\documentclass{rvcepaper}
\begin{document}

% ---------- HS361TA Entrepreneurship & IPR, Test-2, May 2026
\begin{iempaperB}
\paperinfo{hdr=2,usn=0,ayear={2025-2026 (Even Sem)},dept={INDUSTRIAL ENGINEERING AND MANAGEMENT},date={MAY 2026},code={HS361TA},sem={VI Semester},exam={Test-2},maxmarks={50},duration={90 Min},
 title={Entrepreneurship and Intellectual Property Rights},shownote=0}
\iemheadB
\ptitle{PART -- A}
\begin{xtable}{Questions}
\q{1}{What is the primary objective of registering an Industrial Design?}{02}{L2}{5}
\q{2}{Copy Right protects expression but not the ideas. Justify?}{02}{L2}{5}
\q{3}{Define Performers Rights.}{02}{L1}{5}
\q{4}{Define `Organic Growth' in the context of business expansion.}{02}{L2}{1}
\q{5}{What is a Break-Even Point in the context of business operations?}{02}{L2}{3}
\end{xtable}
\ptitle{PART -- B}
\begin{xtable}{Questions}
\q{1.}{A company plans to launch an electric scooter in a semi-urban market where fuel prices are rising, but charging infrastructure is limited. Explain how the company should assess market feasibility and demand before launching the product.}{10}{L3}{2}
\q{2.}{A startup is launching an AI-based online tutoring platform connecting freelance tutors with students across India. Explain what key elements should be included in the executive summary of its business plan.}{10}{L4}{3}
\q{3.}{A food startup has developed a unique recipe and process for manufacturing a long-lasting organic juice. The company wants to prevent competitors from copying it. Explain how the firm can protect its trade secrets in India and why such protection is significant.}{10}{L3}{5}
\q{4.}{An educator creates a comprehensive online course including recorded lectures, PDFs, animations, and quizzes, and uploads it on a learning platform. Later, another platform copies the videos and distributes them without permission. Explain the nature, scope, and rights violated in this situation.}{10}{L3}{5}
\q{5.}{A startup proposes an AI-based mobile app that detects plant diseases using images. Explain how the startup should evaluate technical feasibility through prototype development and proof of concept before full-scale launch.}{10}{L3}{4}
\end{xtable}
\btfoot{|l|l|l|*{11}{C{0.8cm}|}}{\hline
\multirow{2}{*}{\shortstack{Marks\\Distribution}} & \multicolumn{2}{c|}{Particulars} & CO1 & CO2 & CO3 & CO4 & CO5 & L1 & L2 & L3 & L4 & L5 & L6\\\cline{2-14}
 & Max & Quiz & 06 & 04 & & & & 04 & 06 & & & & \\\cline{3-14}
 & Marks & Test & 10 & 10 & 20 & 10 & & & 30 & 10 & 10 & & \\\hline}
\end{iempaperB}

% ---------- HS361TA Entrepreneurship & IPR, CIE - II, 6th Apr 2025
\begin{iempaper}
\paperinfo{shownote=0,date={6\textsuperscript{th} Apr 2025},exam={CIE -- II},maxmarks={10 + 50},sem={VI},prog={UG},duration={30 + 90 Min},
 title={Entrepreneurship \& Intellectual Property Rights},code={HS361TA}}
\iemhead
{\small Note:}\par\vspace{-1pt}
\hspace*{0.5cm}{\footnotesize 1. \hspace{0.1cm}Answer all the Questions.}\par\vspace{3pt}
\begin{cietable}{M}{BT}{CO}
\qpart{Part -- A}
\q{1.}{Define Eco-Labelling.}{2}{2}{4}
\q{2.}{Differentiate between trademark and passing off. (Any two points)}{2}{2}{5}
\q{3.}{What is a Collective Trademark?}{2}{2}{4}
\q{4.}{What is a Break-even point in the context of business operations?}{2}{2}{3}
\q{5.}{Why is a financial plan important in a business plan?}{2}{2}{3}
\qpart{Part -- B}
\q{1.}{Discuss any two important trademark infringement cases, highlighting their impact on trademark law.}{10}{3}{3}
\q{2.}{Explain the difference between registrable and non-registrable trademarks. Discuss the criteria that make a trademark registrable and provide examples of marks that cannot be registered.}{10}{3}{4}
\q{3a.}{The Soil Association, a UK-based organic certification organization, owned a registered eco-label that certified products meeting strict organic farming standards. The Association of British Growers (ABG) used a mark similar to the Soil Association's eco-label on their products without permission. Whether ABG's unauthorized use of the Soil Association's eco-label constituted trademark infringement and whether it misled consumers into believing the products were certified organic by the Soil Association. Discuss.}{05}{3}{5}
\q{3b.}{ITC Ltd., an Indian conglomerate, owned the trademark ``Munch'' for its chocolate products. Nestle India Ltd. started marketing a chocolate bar also named ``Munch.'' ITC alleged that Nestle's use of the mark was deceptively similar and infringed its trademark rights.
\begin{bul}\item Whether Nestl\'e's use of the mark ``Munch'' infringed ITC's registered trademark by creating deceptive similarity.\item Whether customers were likely to be confused or misled by the similarity of the marks.\item Is Nestl\'e's use amounted to unfair competition or passing off.\end{bul}}{05}{4}{5}
\q{4.}{Discuss the following aspects related to new product development and innovation:\newline
a) How can identifying market opportunities and trends influence the ideation process?\newline
b) Explain the role of integration of engineering principles and cross-disciplinary collaboration in technological innovation.\newline
c) Describe the methods to assess market feasibility and demand analysis for a new product.\newline
d) Discuss how technical feasibility is evaluated through prototype development and proof of concept.\newline
e) Outline the key components of financial feasibility analysis including cost estimation and revenue projection.}{10}{2}{3}
\q{5.}{Describe the SWOC analysis and explain Porter's generic competitive strategies with examples.}{10}{3}{3}
\end{cietable}
\end{iempaper}

\end{document}
